% Figure: sensitivity of the double pendulum. Two trajectories with
% initial conditions differing by 1e-3 rad diverge after a few seconds;
% the angle difference grows roughly exponentially then saturates.
\begin{figure}[htbp]
\centering
\begin{tikzpicture}[
  axis/.style={draw=engblack, -{Stealth}, thick},
  growth/.style={draw=engred, very thick},
  envel/.style={draw=engblue, very thick, dashed},
  guide/.style={draw=enggray, dashed, thin},
  label/.style={font=\small},
]
% x = time 0..12 s (plotted 0..10), y = log10 of angle difference, -6..1 (plotted 0..5)
\draw[axis] (0,0) -- (10.6,0) node[anchor=west, label] {$t$ [s]};
\draw[axis] (0,0) -- (0,5.6)  node[anchor=south, label] {$\log_{10}|\Delta\theta_2|$};

% x ticks (1.2 s per unit)
\foreach \x/\m in {0/0, 2.5/3, 5/6, 7.5/9, 10/12} {
  \draw (\x,-0.08) -- (\x,0.08);
  \node[font=\scriptsize, below] at (\x,-0.12) {\m};
}
% y ticks: log10 from -6 (y=0) to +1 (y=5); step 1.4 per decade
\foreach \y/\h in {0/$-6$, 1.43/$-4$, 2.86/$-2$, 4.29/$0$} {
  \draw (-0.08,\y) -- (0.08,\y);
  \node[font=\scriptsize, left] at (-0.12,\y) {\h};
}

% Linear-in-log growth = exponential divergence, slope = Lyapunov exponent
% start at log10 = -6 (separation 1e-6 in theta1 maps to ~1e-6 in theta2 after coupling)
% reach saturation ~ log10 = 0.3 (order 1 rad) by t ~ 9 s, then flatten
% slope: from (0,0) to (7.5, 4.3): straight line
\draw[growth] (0,0) -- (7.0, 4.3);
% saturation plateau
\draw[growth] (7.0,4.3) .. controls (8.0,4.6) and (9.0,4.55) .. (10,4.5);
\node[engred, label, anchor=south west] at (2.6, 1.6) {$|\Delta\theta_2|\sim e^{\lambda t}$};

% slope guide annotation
\draw[guide] (1,0.61) -- (3,0.61) -- (3, {0.61 + 2*0.61});
\node[font=\scriptsize, enggray, anchor=west] at (3.1, 1.0) {slope $=\lambda/\ln 10$};

% saturation note
\node[engblue, font=\scriptsize, anchor=south east] at (10, 4.8) {saturation $\sim O(1)$ rad};
\draw[guide] (0, 4.5) -- (10, 4.5);

\end{tikzpicture}
\caption[Exponential divergence of nearby double-pendulum
trajectories]{The signature of deterministic chaos in the double
pendulum. Two simulations start from initial conditions differing by
$10^{-6}\,\si{\radian}$ in the upper angle. The base-ten logarithm of
the resulting difference $|\Delta\theta_2(t)|$ grows linearly with time
(red), which means the difference itself grows exponentially,
$|\Delta\theta_2|\sim e^{\lambda t}$; the slope divided by $\ln 10$ is
the largest Lyapunov exponent $\lambda$. Growth continues until the two
trajectories are fully uncorrelated, where the separation saturates at
order one radian (blue dashed). The energy method delivers the exact
equations of motion, and those exact equations are unpredictable in
practice: determinism and predictability are not the same property.}
\label{fig:vol03ch05:double-pendulum-divergence}
\end{figure}
